\section{The Architect's Question}\label{the-architects-question}

After this chapter we should be able to treat AI delivery as an
\emph{industrial process} rather than a one-off project: a repeatable
factory with automated pipelines for data, evaluation, promotion, and
rollback. We will separate the three layers of the factory (the model
lifecycle, the serving pipeline, and the operations loop), then build
the economics of promotion gating and retraining cadence for the
canonical workload. After this chapter, ``we ship a better model'' is a
gated, measurable, reversible event --- not a risky redeploy.

\section{Concept}\label{concept}

The AI Factory is the discipline of making model delivery routine. It
rests on three pillars:

\begin{enumerate}
\def\labelenumi{\arabic{enumi}.}
\item
  \textbf{Data pipeline} --- the automated intake, cleaning, versioning,
  and labelling of the data that trains or evaluates the model. For a
  RAG system this includes the retrieval corpus, its freshness, and its
  quality gates. {[}2° FACT{]}
\item
  \textbf{Evaluation \& promotion gate} --- a systematic process that
  scores a candidate model or configuration against defined metrics
  (quality, latency, cost) on a held-out benchmark before it can replace
  the incumbent in production. This is where Ch14's deployment
  benchmarking becomes a recurring gate rather than a one-time event.
\item
  \textbf{Serve + observe + rollback} --- the serving layer (Ch11) plus
  a feedback loop that measures real traffic, detects drift or
  regression, and can roll the production model back safely. {[}2°
  FACT{]}
\end{enumerate}

The factory converts the architect's recurring question from ``should we
adopt this model?'' (a one-off decision) into ``does this candidate pass
the automated gate?'' (a routine, repeatable decision with the same
rigor every time).

\section{Mental Model}\label{mental-model}

Think of the AI Factory as a \textbf{release pipeline with a quality
gate placed before every promotion}, mirroring how software CI/CD gated
deployments decades ago. Three mental moves:

\begin{itemize}
\tightlist
\item
  \textbf{Models are artifacts like code.} They get versioned, tested,
  reviewed, and promoted through environments (staging → canary →
  production) with automated checks at each step --- not dropped into
  prod on a hunch.
\item
  \textbf{The gate is the memory of the factory.} The evaluation suite
  and its thresholds encode everything we've learned about what ``good
  enough'' means (Ch4 capability, Ch14 deployment). If the gate is weak,
  the factory has no standards.
\item
  \textbf{Everything is reversible.} A promotion that regresses metrics
  must be rollback-able in minutes. The factory's value is not
  eliminating errors but making them cheap and safe.
\end{itemize}

The factory's shape is a \textbf{pipeline with decisions}: data in →
train/fine-tune → evaluate → push to canary → observe → promote or
rollback → serve. Each arrow is a decision point, and each decision
point is where an architect has influence.

\section{Worked Example}\label{worked-example}

\subsection{The Economics of Promotion Gating for the Canonical
Workload}\label{the-economics-of-promotion-gating-for-the-canonical-workload}

We show the factory's value by comparing \emph{ungated} and \emph{gated}
adoption of a candidate model on the canonical 70B RAG workload
(\textasciitilde10 rps average, \textsubscript{9.2K-in/}300-out). {[}1P
§14{]}

\textbf{Setup.} Suppose a new model candidate claims better retrieval-QA
quality but no one has measured its serving behavior. Two adoption
paths:

\begin{itemize}
\tightlist
\item
  \textbf{Ungated}: deploy the candidate to all production traffic
  immediately. If it regresses (say p95 TTFT degrades because of a
  slower tokenizer, or quality drops on a subset), the service suffers
  until rollback.
\item
  \textbf{Gated}: run the candidate on a 5\% canary for N hours, measure
  TTFT/goodput/quality against the incumbent, promote only if it beats
  every threshold.
\end{itemize}

\textbf{Cost of a bad ungated release.} The canonical service does
\textasciitilde864,000 requests/day. Suppose a regression raises p95
TTFT from 1.9 s to 2.8 s (breaching the 2 s SLO) for 6 hours before
discovery and rollback. That is 6/24 × 864,000 ≈ \textbf{216,000
requests} served at a broken SLO. At \textasciitilde\$0.82/1K requests
(Ch16) the direct cost is modest (\textasciitilde\$180), but the
\emph{business} cost --- users experiencing a 47\% latency regression,
trust erosion, and the fire drill --- is the dominant term. {[}2°
DERIVED{]} \emph{(worked example, not reference)}

\textbf{Cost of a good gated release.} A canary at 5\% traffic for 2
hours serves 5\% × 864,000 × 2/24 ≈ \textbf{3,600 requests} through the
candidate while it is measured against a full evaluation suite. If it
passes, promotion is low-risk; if it fails, no users beyond canary were
exposed. {[}2° DERIVED{]}

\textbf{Reading the result.} Gating turns a risky 216K-request exposure
(ungated) into a controlled \textasciitilde3.6K-request canary --- a
\textasciitilde60× reduction in blast radius for the cost of a 2-hour
evaluation. The factory makes this routine. {[}2° DERIVED{]}

\textbf{\hyperref[tab:21.1]{Table 21-1}} --- Ungated vs gated model adoption (illustrative
worked example)

\begin{longtable}[]{@{}
  >{\raggedright\arraybackslash}p{0.1667\linewidth}
  >{\raggedright\arraybackslash}p{0.1667\linewidth}
  >{\raggedright\arraybackslash}p{0.1667\linewidth}
  >{\raggedright\arraybackslash}p{0.1667\linewidth}
  >{\raggedright\arraybackslash}p{0.1667\linewidth}
  >{\raggedright\arraybackslash}p{0.1667\linewidth}@{}}
\toprule\noalign{}
\begin{minipage}[b]{\linewidth}\raggedright
Adoption
\end{minipage} & \begin{minipage}[b]{\linewidth}\raggedright
Requests exposed
\end{minipage} & \begin{minipage}[b]{\linewidth}\raggedright
SLO breach risk
\end{minipage} & \begin{minipage}[b]{\linewidth}\raggedright
Discovery
\end{minipage} & \begin{minipage}[b]{\linewidth}\raggedright
Rollback
\end{minipage} & \begin{minipage}[b]{\linewidth}\raggedright
Best when
\end{minipage} \\
\midrule\noalign{}
\endhead
\bottomrule\noalign{}
\endlastfoot
Ungated & \textasciitilde216K in a bad release & high (all traffic) &
post-hoc & slow & tiny/experimental \\
Gated (5\% canary) & \textasciitilde3.6K & low (confined) & automated &
fast & business-critical serving \\

\label{tab:21.1}\end{longtable}

\emph{(Worked-example numbers for a 2-hour canary at 5\% traffic; {[}2°
DERIVED{]} illustrative, {[}VERIFY{]} per deployment.)}

\section{Measurement}\label{measurement}

The factory is only as trustworthy as its gates. We measure four things
continuously:

\begin{enumerate}
\def\labelenumi{\arabic{enumi}.}
\tightlist
\item
  \textbf{Gate pass rate / false-promotion rate} --- how often a gated
  change later regresses in production; a rising value means the
  evaluation suite is losing fidelity.
\item
  \textbf{Canary-to-production time} and \textbf{rollback time} --- the
  speed of promotion and escape; both should be minutes, not days.
\item
  \textbf{Model freshness / data staleness} --- time since the retrieval
  corpus or model was last validated; drift detection feeds the
  retraining decision.
\item
  \textbf{MTTR for a bad release} --- mean time to detect + rollback;
  the factory's core reliability number.
\end{enumerate}

\section{Common Mistakes}\label{common-mistakes}

\begin{itemize}
\tightlist
\item
  \textbf{No gate, or a decorative gate.} An evaluation that always
  passes is worse than none because it grants false confidence.
\item
  \textbf{Promoting on quality only.} Quality passing but serving (Ch14)
  regressing still breaks the service; the gate must include deployment
  metrics.
\item
  \textbf{Assuming retraining cadence is ``as often as possible.''} More
  frequent retraining is not free --- it consumes compute, risk, and
  evaluation effort; the right cadence balances drift against churn (see
  Ch7-8 for training economics).
\item
  \textbf{Skipping rollback rehearsal.} A rollback that has never been
  exercised will fail precisely when it is needed.
\item
  \textbf{Treating the factory as a build-once artifact.} It is a living
  process that must absorb new models, new metrics, and new data sources
  continuously.
\end{itemize}

\section{Architecture Consequence}\label{architecture-consequence}

The AI Factory is the operational backbone that makes the book's whole
loop durable. Every architectural decision recomputes through it: model
selection (Ch3/5) feeds the evaluation gate, serving choices (Ch11)
determine canary mechanics, benchmarking (Ch14) becomes the recurring
gate, TCO (Ch16) sets how much gating overhead is affordable, and fleet
operations (Ch17-20) provide the rollout/rollback surfaces. Without the
factory, each chapter's decision is a one-off; with it, they compound
into a system that gets better and safer every cycle.

\begin{figure}
\centering
\pandocbounded{\includegraphics[keepaspectratio,alt={\hyperref[fig:21.1]{Fig 21.1} --- The AI Factory promotion pipeline {[}ILLUSTRATIVE conceptual{]}},width=\maxwidth{\textwidth},keepaspectratio]{figures/fig-21-2101.pdf}}
\caption{\hyperref[fig:21.1]{Fig 21.1} --- The AI Factory promotion pipeline {[}ILLUSTRATIVE
conceptual{]}}

\label{fig:21.1}\end{figure}

\emph{\hyperref[fig:21.1]{Fig 21.1} --- The model lifecycle: data → evaluate gate → canary →
promote/rollback, with an observation feedback loop.}

\section{What We Still Don't Know}\label{what-we-still-dont-know}

\begin{itemize}
\tightlist
\item
  \textbf{The true cost of model churn} (retraining + reeval + redeploy)
  per productivity gain is {[}HYPOTHESIS{]} and workload-specific; there
  is no universal cadence.
\item
  \textbf{How well evaluation suites generalize} across distribution
  shifts in real traffic is {[}VERIFY{]}; a suite tuned on last
  quarter's data may misjudge this quarter's model.
\item
  \textbf{The right canary size} for a given risk posture is
  {[}VERIFY{]} per service; 5\% worked here but high-risk systems may
  need smaller, staged canaries.
\end{itemize}

\section{End-of-Chapter Mini-Case}\label{end-of-chapter-mini-case}

A vendor pitches the architect a newer, cheaper model: ``just swap it in
--- it scores better on the bench.'' A year ago the architect might
have, and watched p95 TTFT silently drift past the SLO for a week. Now
the AI Factory is in place. The architect drops the candidate into the
pipeline: the data pipeline versions it, the evaluation gate scores it
against quality AND the deployment benchmark (TTFT/goodput/cost) on the
canonical workload, and a 5\% canary serves it for two hours. The gate
flags a quality regression on the legal-document subset the benchmark
suite specifically covers --- a subset the vendor's ``better bench''
never measured. The candidate is rejected automatically, the incumbent
stays, and the legal subset keeps working. No fire drill, no
user-visible regression, no manual heroics --- the factory made the
wrong vendor claim cost nothing to the business, which is precisely what
industrializing the delivery loop is for.
